\section{Prime faces and the complete Weil functional}
\label{sec:prime-faces-weil}
\subsection{The exact semilocalization associated with a monoid prime}
For $\Sigma\subseteq\mathcal P$, define
\[
 T_\Sigma=\{n\in\Z\setminus\{0\}:p\nmid n\text{ for every }p\in\Sigma\},
 \qquad R_\Sigma=T_\Sigma^{-1}\Z.
\]
\begin{theorem}\label{thm:face-semilocalization}
The complement of the monoid prime $J_\Sigma$ is $T_\Sigma$. Localization followed by the original supported addition gives exactly $G(R_\Sigma)$. Its arithmetic spectrum consists of $(0)$ and the primes $p\in\Sigma$, each with its original residue field. If $F\subset\mathcal P$ is finite, then
\begin{equation}\label{eq:face-S-units}
 R_{\mathcal P\setminus F}=\Z[p^{-1}:p\in F],\qquad
 R_{\mathcal P\setminus F}^{\times}
 =\{\pm\prod_{p\in F}p^{n_p}:n_p\in\Z\}=\Gamma_F.
\end{equation}
\end{theorem}
\begin{proof}
The complement statement follows directly from the definition of $J_\Sigma$. A supported fraction $a/t$ obeys the original fraction operations and $\tau/t=\tau$, giving the localization isomorphism in Theorem~\ref{thm:rebuilt-scalar}. A ring prime survives localization exactly when it misses $T_\Sigma$: $(p)$ misses it precisely when $p\in\Sigma$, and $(0)$ always misses it. Its localized residue is still $\F_p$, since the denominators are nonzero modulo $p$. For the complement of $F$, the allowed denominators have no prime factors outside $F$, yielding the first equality in \eqref{eq:face-S-units}. A rational number and its inverse belong to that localization exactly when numerator and denominator have no prime factors outside $F$, proving the unit formula.
\end{proof}
The arithmetic trace space associated with these units is
\[
 \mathbb A_F=\R\times\prod_{p\in F}\Q_p,\qquad
 Y_F=\mathbb A_F/\Gamma_F.
\]
This is the precise set-of-places and scaling-group construction used by Connes \cite[labels \texttt{YQS}, \texttt{GL1QS}]{Connes2026}. Equation~\eqref{eq:face-S-units} identifies its group with the original multiplicative-prime localization, without replacing that spectrum by the scalar semiring spectrum.

\subsection{Every prime-power, Gamma and pole term}
Use the multiplicative convention
\[
 \widehat h(s)=\int_0^\infty h(x)x^{-is}\frac{dx}{x},\quad
 g^*(x)=\overline{g(x^{-1})},\quad
 (g_1*g_2)(x)=\int_0^\infty g_1(y)g_2(x/y)\frac{dy}{y}.
\]
For $h\in C_c^\infty(\R_+^*)$, set
\begin{align}
 W_p(h)&=(\log p)\sum_{m\ge1}p^{-m/2}\bigl(h(p^m)+h(p^{-m})\bigr),\label{eq:rebuilt-Wp}\\
 W_\infty(h)&=(\log(4\pi)+\gamma_E)h(1)\notag\\
 &\quad+\int_1^\infty\left(h(x)+h(x^{-1})-2x^{-1/2}h(1)\right)
 \frac{x^{1/2}}{x-x^{-1}}\frac{dx}{x},\label{eq:rebuilt-Winf}\\
 D_F(h)&=\widehat h(i/2)+\widehat h(-i/2)-W_\infty(h)-\sum_{p\in F}W_p(h).
 \label{eq:rebuilt-D}
\end{align}
The cancellation in the integrand at $x=1$ is retained; beyond the support of $h$, its subtraction term is also retained. These are the explicit-formula conventions in \cite[\texttt{wp}, \texttt{arch}, \texttt{sectweilpos}]{Connes2026}. In this notation $W_\infty$ is that source's $W_{\mathbb R}$, without changing its sign. The full Weil functional is $D_{\mathcal P}$, whose prime sum is finite for each such $h$.

\begin{theorem}[The entire prime-face valuation]\label{thm:prime-face-weil}
For finite sets of rational primes $F,E$,
\begin{align}
 D_{F\cup E}(h)+D_{F\cap E}(h)&=D_F(h)+D_E(h),\label{eq:face-valuation}\\
 D_{F\cup\{p\}}(h)-D_F(h)&=-W_p(h)\quad(p\notin F).
 \label{eq:face-difference}
\end{align}
If $\supp g\subseteq[\lambda^{-1},\lambda]$ and $h=g*g^*$, every prime term with $p>\lambda^2$ vanishes. Thus for $F_\lambda=\{p:p\le\lambda^2\}$,
\begin{equation}\label{eq:full-Weil-return}
 D_{F_\lambda}(g*g^*)=D_{\mathcal P}(g*g^*).
\end{equation}
The full Weil positivity criterion therefore concerns these compatible complete finite-place returns. Equations \eqref{eq:face-valuation}--\eqref{eq:face-difference} specify exactly what changes when a prime face is omitted.
\end{theorem}
\begin{proof}
In \eqref{eq:rebuilt-D}, both pole terms and the archimedean term are independent of $F$. Each prime occurs with the indicator $1_F(p)$, and
$1_{F\cup E}+1_{F\cap E}=1_F+1_E$, proving the first formula. Subtraction gives the second with its minus sign. In a nonzero convolution integrand for $h(x)$, both $y$ and $x/y$ belong to $[\lambda^{-1},\lambda]$, since the involution preserves that interval. Hence $x\in[\lambda^{-2},\lambda^2]$. For $p>\lambda^2$ and $m\ge1$, neither $p^m$ nor $p^{-m}$ belongs to that interval, so every summand in \eqref{eq:rebuilt-Wp} is zero. This proves \eqref{eq:full-Weil-return}. The classical equivalence of positivity of this full functional and RH is Weil's criterion as stated in the cited source; the finite-face identities themselves have been proved here without assuming RH.
\end{proof}

\begin{theorem}[Prime-face increments have both signs, even with both poles killed]\label{thm:prime-face-signs}
Fix a rational prime $p$, put $a=\log p$, and choose a nonzero real $\phi\in C_c^\infty((-\eta,\eta))$ with $0<4\eta<a$. Set
\[
 \psi=\phi''-\tfrac14\phi,\qquad
 f_\theta(u)=\psi(u)+e^{i\theta}\psi(u-a),\qquad
 g_\theta(x)=f_\theta(\log x).
\]
Then $\widehat g_\theta(i/2)=\widehat g_\theta(-i/2)=0$, and for $h_\theta=g_\theta*g_\theta^*$,
\begin{equation}\label{eq:prime-increment-evaluated}
 W_p(h_\theta)=2(\log p)p^{-1/2}\cos\theta\,\|\psi\|_{L^2(\R)}^2.
\end{equation}
In particular the change \eqref{eq:face-difference} is strictly negative at $\theta=0$ and strictly positive at $\theta=\pi$. Retaining residue norm $p$ does not make prime-face removal a sign-neutral operation.
\end{theorem}
\begin{proof}
If $\psi=0$, then $\phi''=\phi/4$. Multiplying by $\phi$ and integrating gives $-\int|\phi'|^2=\frac14\int|\phi|^2$, forcing $\phi=0$, a contradiction. Thus $\|\psi\|^2>0$.
Integration by parts twice gives
$\int\psi(u)e^{\pm u/2}\,du=0$. Translation multiplies these integrals by $e^{\pm a/2}$, so both Mellin values remain zero for $g_\theta$.

In logarithmic coordinates, $h_\theta(e^v)=\int f_\theta(u)\overline{f_\theta(u-v)}\,du$. At $v=a$, only the right bump of the first factor and the left bump of the second overlap; the other three products vanish since $2\eta<a$. Their integral is $e^{i\theta}\|\psi\|^2$. At $v=-a$ it is its conjugate. For $|m|\ge2$, the shifts $ma$ exceed the difference support $[-a-2\eta,a+2\eta]$, so the correlations are zero. Substitution in the full sum \eqref{eq:rebuilt-Wp} proves \eqref{eq:prime-increment-evaluated}. The pole conditions hold exactly, without deleting either pole from the definition of $D_F$.
\end{proof}

This calculation is a signed arithmetic consequence of the prime-face map. It does not assign positivity to the complete functional: the other prime terms and the archimedean term in \eqref{eq:rebuilt-D} remain part of that calculation. It supplies an explicit reason that no argument may infer the full Weil sign from an unchanged list of residue cardinalities.
